\section{引言}
\label{sec:introduction}

边缘目标检测要求在有限功耗和存储资源下持续处理单张图像。卷积提供了网络中
的大部分乘加运算，因此固定规模脉动阵列成为常用实现方式
\cite{kung1982systolic,zhang2015optimizing,qiu2016going}。然而，网络层的规约
维长度通常远大于阵列行数。硬件必须把一个卷积拆成多个 $K$ 维分段计算轮次，
并把输出通道拆成多个并行块。阵列规模决定峰值运算能力，分段后的数据供应方式
则决定这些运算单元能否连续工作。

完整 YOLOv3-tiny 是这一矛盾的代表性工作负载\cite{yolov3}。在固定
416$\times$416 输入下，13 个卷积的规约长度从 27 变化到 4608，空间尺寸从
416$\times$416 变化到 13$\times$13，输出通道还包含不能被 32 整除的两个
255 通道检测头。对本文的 18 行阵列，浅层 m0 只需要 2 个计算轮次，深层 m13
则需要 256 个计算轮次和 32 个输出通道块。若每个计算上下文重新从 DDR 获取
同一输入 tile，输入流量会随两维分块乘积放大；即使软件预先生成卷积窗口，
展开数据仍需按输出通道块重复传输。由此形成本文关注的单一核心冲突：
\textbf{长 $K$ 卷积需要大量分段复用，而串行逐轮流送把这种复用转化为重复
片外搬运和计算边界空泡。}

现有 CNN 加速器通过循环分块、数据流映射和可重构互连提高数据复用
\cite{chen2016eyeriss,eyerissv2,maeri2018,simba2019}，FPGA 工具链则从网络
描述生成定制流水或折叠阵列\cite{fpgaConvNet2015,finn2017,dnnBuilder2018}。
这些研究奠定了片上复用基础。面向软件可见 HWC 张量和多轮固定阵列，仍需把
“输入只进入 PL 一次”落实为可在反压、尾块和分支网络下执行的硬件语义。
LASA 的出发点正是将输入布局转换和跨轮复用合并为一个片上过程：原始 HWC
tile 进入窗口生成单元后形成阵列向量，后续所有 $K$ 维轮次和输出通道块直接
访问片上缓存。权重提前装载、部分和反馈和上下文切换随后围绕这一数据流进行
调度，使缓存复用能够转化为连续计算。

本文提出 \LASA（Layer-Adaptive Systolic Accelerator）。其 18 路规约输入与
16 列双输出通道结构在一个周期完成 576 次 INT8 乘加，并以 32 通道为一个输出
块。3$\times$3 卷积从 HWC 输入生成窗口并物化，1$\times$1 卷积使用直接向量
路径；两种路径共享阵列、部分和反馈、重量化和 HWC 写回。空间 tile 高度根据
层形状和片上容量选择，使浅层降低边界启动次数、长层满足窗口与部分和容量。

本文的主要贡献概括为以下两点。

\begin{enumerate}[leftmargin=2.1em]
  \item \textbf{片上物化重放的数据流。}LASA 将软件可见的 HWC tile 一次转换
  为阵列窗口，并跨 $K$ 维计算轮次和输出通道块重放。本文建立片外输入流量、
  片上物化流量和计算时延模型，说明该机制如何解除 DDR 输入量与计算上下文数
  之间的耦合；层自适应 tile、原生 1$\times$1/3$\times$3 路径和尾通道处理
  共同保证完整网络能够映射到同一阵列。
  \item \textbf{面向连续计算的上下文调度。}LASA 在保持缓冲区唯一占用、数据
  不覆盖和顺序提交的条件下，提前装载下一组权重，并把部分和反馈、结果提交和
  下一上下文准备与当前计算重叠，从而缩短物化重放数据流中的轮次边界。
\end{enumerate}

本文在 KV260/XCK26 上实现完整双尺度 YOLOv3-tiny。全部 13 个卷积在 PL
执行，池化、上采样、分支重量化拼接和双尺度后处理由四核 Cortex-A53 完成。
200 MHz 实现取得 \ResidentFPS{} FPS、\PLEffectiveGOPS{} GOPS 和
\ArrayUtilization{}\% 阵列利用率；\ConformanceImages{} 张图的
\NodeRecords{} 条整数节点记录全部逐字节一致，5000 张 COCO val2017 的原始
检测头指标与主机参考完全相同。这些结果把理论流量模型、硬件实现和完整网络
行为统一到同一推理系统中。

本文其余部分组织如下：\cref{sec:problem} 建立工作负载和理论模型；
\cref{sec:architecture} 介绍 LASA 体系结构；\cref{sec:materialize} 说明片上
物化重放和上下文调度；\cref{sec:system-method} 给出完整网络协同和实验方法；
\cref{sec:results} 报告实验结果；最后总结适用范围和结论。
